\documentclass[12pt, a4paper]{article}

\usepackage[utf8]{inputenc}
\usepackage[mathscr]{euscript}
\let\euscr\mathscr \let\mathscr\relax
\usepackage[scr]{rsfso}
\usepackage{amssymb,amsmath,amsthm,amsfonts}
\usepackage[shortlabels]{enumitem}
\usepackage{multicol,multirow}
\usepackage{lipsum}
\usepackage{balance}
\usepackage{calc}
\usepackage[colorlinks=true,citecolor=blue,linkcolor=blue]{hyperref}
\usepackage{import}
\usepackage{xifthen}
\usepackage{pdfpages}
\usepackage{transparent}
\usepackage{tabularx}
\usepackage{listings}
\usepackage[dvipsnames]{xcolor}
\usepackage{bbm}

\newcommand{\incfig}[2][1.0]{
    \def\svgwidth{#1\columnwidth}
    \import{./figures/}{#2.pdf_tex}
}
\newcommand{\incimg}[2][1.0]{
  \includegraphics[width=#1\columnwidth]{./figures/#2}
}

\newcommand{\xor}{\oplus}
\newcommand{\floor}[1]{\left\lfloor #1 \right\rfloor}
\newcommand{\ceil}[1]{\left\lceil #1 \right\rceil}

\newlist{enumproof}{enumerate}{4}
\setlist[enumproof,1]{label=\arabic*., parsep=1em}
\setlist[enumproof,2]{label=\arabic{enumproofi}.\arabic*., parsep=1em}
\setlist[enumproof,3]{label=\arabic{enumproofi}.\arabic{enumproofii}.\arabic*., parsep=1em}
\setlist[enumproof,4]{label=\arabic{enumproofi}.\arabic{enumproofii}.\arabic{enumproofiii}.\arabic*., parsep=1em}

\lstdefinestyle{mystyle}{
  language=Python, % Set the language to Python
  commentstyle=\color{codegray}\itshape, % Python comments (#)
  keywordstyle=\color{blue},  % Python reserved keywords
  stringstyle=\color{mauve},             % Strings
  basicstyle={\ttfamily\footnotesize},   % Monospace font
  breakatwhitespace=true,         
  breaklines=true,                 
  captionpos=b,                    
  keepspaces=true,                 
  numbers=none,                    
  tabsize=2,
  showstringspaces=false,
  % Extra keyword groups for Python
  morekeywords=[2]{self, cls, __init__, __str__, __repr__}, 
  keywordstyle=[2]\color{identblue},     % Special identifiers / methods
  morekeywords=[3]{int, float, bool, str, list, dict, set, tuple, object, None, True, False},
  keywordstyle=[3]\color{codegreen},     % Built-in types
  morekeywords=[4]{print, input, open, range, len, type, isinstance, super, enumerate, zip, map, filter, sum, min, max, sorted},
  keywordstyle=[4]\color{cyan},          % Built-in functions
}
\lstset{style=mystyle}

\usepackage{ifthen}
\usepackage[landscape]{geometry}
\usepackage[shortlabels]{enumitem}

\ifthenelse{\lengthtest { \paperwidth = 11in}}
    { \geometry{top=.5in,left=.5in,right=.5in,bottom=.5in} }
	{\ifthenelse{ \lengthtest{ \paperwidth = 297mm}}
		{\geometry{top=1cm,left=1cm,right=1cm,bottom=1cm} }
		{\geometry{top=1cm,left=1cm,right=1cm,bottom=1cm} }
	}

\pagestyle{empty}
\makeatletter
\renewcommand\thesection{\arabic{section}.}
\renewcommand{\section}{\@startsection{section}{1}{0mm}%
                                {-1ex plus -.5ex minus -.2ex}%
                                {0.05ex}%x
                                {\normalfont\normalsize\bfseries}}
\renewcommand{\subsection}{\@startsection{subsection}{2}{0mm}%
                                {-1ex plus -.5ex minus -.2ex}%
                                {0.05ex}%
                                {\normalfont\small\bfseries}}
\renewcommand{\subsubsection}{\@startsection{subsubsection}{3}{0mm}%
                                {-1ex plus -.5ex minus -.2ex}%
                                {0.05ex}%
                                {\normalfont\footnotesize\bfseries}}
\newcommand{\colbreak}{\vfill\null\columnbreak}
\makeatother
\setcounter{secnumdepth}{1}
\setlength{\parindent}{0pt}
\setlength{\parskip}{0.7em}

\setlist[itemize]{itemsep=0.6ex, topsep=-2pt, partopsep=0pt, parsep=0pt}
\setlist[enumerate]{itemsep=0.6ex, topsep=-2pt, partopsep=0pt, parsep=0pt}

% Things Lie
\newcommand{\kb}{\mathfrak b}
\newcommand{\kg}{\mathfrak g}
\newcommand{\kh}{\mathfrak h}
\newcommand{\kn}{\mathfrak n}
\newcommand{\ku}{\mathfrak u}
\newcommand{\kz}{\mathfrak z}
\DeclareMathOperator{\Ext}{Ext} % Ext functor
\DeclareMathOperator{\Tor}{Tor} % Tor functor
\newcommand{\gl}{\opname{\mathfrak{gl}}} % frak gl group
\renewcommand{\sl}{\opname{\mathfrak{sl}}} % frak sl group chktex 6

% More script letters etc.
\newcommand{\SA}{\mathcal A}
\newcommand{\SB}{\mathcal B}
\newcommand{\SC}{\mathcal C}
\newcommand{\SF}{\mathcal F}
\newcommand{\SG}{\mathcal G}
\newcommand{\SH}{\mathcal H}
\newcommand{\OO}{\mathcal O}

\newcommand{\SCA}{\mathscr A}
\newcommand{\SCB}{\mathscr B}
\newcommand{\SCC}{\mathscr C}
\newcommand{\SCD}{\mathscr D}
\newcommand{\SCE}{\mathscr E}
\newcommand{\SCF}{\mathscr F}
\newcommand{\SCG}{\mathscr G}
\newcommand{\SCH}{\mathscr H}

% Mathfrak primes
\newcommand{\km}{\mathfrak m}
\newcommand{\kp}{\mathfrak p}
\newcommand{\kq}{\mathfrak q}

% number sets
\newcommand{\RR}[1][]{\ensuremath{\ifstrempty{#1}{\mathbb{R}}{\mathbb{R}^{#1}}}}
\newcommand{\NN}[1][]{\ensuremath{\ifstrempty{#1}{\mathbb{N}}{\mathbb{N}^{#1}}}}
\newcommand{\ZZ}[1][]{\ensuremath{\ifstrempty{#1}{\mathbb{Z}}{\mathbb{Z}^{#1}}}}
\newcommand{\QQ}[1][]{\ensuremath{\ifstrempty{#1}{\mathbb{Q}}{\mathbb{Q}^{#1}}}}
\newcommand{\CC}[1][]{\ensuremath{\ifstrempty{#1}{\mathbb{C}}{\mathbb{C}^{#1}}}}
\newcommand{\PP}[1][]{\ensuremath{\ifstrempty{#1}{\mathbb{P}}{\mathbb{P}^{#1}}}}
\newcommand{\HH}[1][]{\ensuremath{\ifstrempty{#1}{\mathbb{H}}{\mathbb{H}^{#1}}}}
\newcommand{\FF}[1][]{\ensuremath{\ifstrempty{#1}{\mathbb{F}}{\mathbb{F}^{#1}}}}
% expected value
\newcommand{\EE}{\ensuremath{\mathbb{E}}}
\newcommand{\charin}{\text{ char }}
\DeclareMathOperator{\sign}{sign}
\DeclareMathOperator{\Aut}{Aut}
\DeclareMathOperator{\Inn}{Inn}
\DeclareMathOperator{\Syl}{Syl}
\DeclareMathOperator{\Gal}{Gal}
\DeclareMathOperator{\GL}{GL} % General linear group
\DeclareMathOperator{\SL}{SL} % Special linear group

%---------------------------------------
% BlackBoard Math Fonts :-
%---------------------------------------

%Captital Letters
\newcommand{\bbA}{\mathbb{A}}	\newcommand{\bbB}{\mathbb{B}}
\newcommand{\bbC}{\mathbb{C}}	\newcommand{\bbD}{\mathbb{D}}
\newcommand{\bbE}{\mathbb{E}}	\newcommand{\bbF}{\mathbb{F}}
\newcommand{\bbG}{\mathbb{G}}	\newcommand{\bbH}{\mathbb{H}}
\newcommand{\bbI}{\mathbb{I}}	\newcommand{\bbJ}{\mathbb{J}}
\newcommand{\bbK}{\mathbb{K}}	\newcommand{\bbL}{\mathbb{L}}
\newcommand{\bbM}{\mathbb{M}}	\newcommand{\bbN}{\mathbb{N}}
\newcommand{\bbO}{\mathbb{O}}	\newcommand{\bbP}{\mathbb{P}}
\newcommand{\bbQ}{\mathbb{Q}}	\newcommand{\bbR}{\mathbb{R}}
\newcommand{\bbS}{\mathbb{S}}	\newcommand{\bbT}{\mathbb{T}}
\newcommand{\bbU}{\mathbb{U}}	\newcommand{\bbV}{\mathbb{V}}
\newcommand{\bbW}{\mathbb{W}}	\newcommand{\bbX}{\mathbb{X}}
\newcommand{\bbY}{\mathbb{Y}}	\newcommand{\bbZ}{\mathbb{Z}}

%---------------------------------------
% MathCal Fonts :-
%---------------------------------------

%Captital Letters
\newcommand{\mcA}{\mathcal{A}}	\newcommand{\mcB}{\mathcal{B}}
\newcommand{\mcC}{\mathcal{C}}	\newcommand{\mcD}{\mathcal{D}}
\newcommand{\mcE}{\mathcal{E}}	\newcommand{\mcF}{\mathcal{F}}
\newcommand{\mcG}{\mathcal{G}}	\newcommand{\mcH}{\mathcal{H}}
\newcommand{\mcI}{\mathcal{I}}	\newcommand{\mcJ}{\mathcal{J}}
\newcommand{\mcK}{\mathcal{K}}	\newcommand{\mcL}{\mathcal{L}}
\newcommand{\mcM}{\mathcal{M}}	\newcommand{\mcN}{\mathcal{N}}
\newcommand{\mcO}{\mathcal{O}}	\newcommand{\mcP}{\mathcal{P}}
\newcommand{\mcQ}{\mathcal{Q}}	\newcommand{\mcR}{\mathcal{R}}
\newcommand{\mcS}{\mathcal{S}}	\newcommand{\mcT}{\mathcal{T}}
\newcommand{\mcU}{\mathcal{U}}	\newcommand{\mcV}{\mathcal{V}}
\newcommand{\mcW}{\mathcal{W}}	\newcommand{\mcX}{\mathcal{X}}
\newcommand{\mcY}{\mathcal{Y}}	\newcommand{\mcZ}{\mathcal{Z}}


%---------------------------------------
% Bold Math Fonts :-
%---------------------------------------

%Captital Letters
\newcommand{\bmA}{\boldsymbol{A}}	\newcommand{\bmB}{\boldsymbol{B}}
\newcommand{\bmC}{\boldsymbol{C}}	\newcommand{\bmD}{\boldsymbol{D}}
\newcommand{\bmE}{\boldsymbol{E}}	\newcommand{\bmF}{\boldsymbol{F}}
\newcommand{\bmG}{\boldsymbol{G}}	\newcommand{\bmH}{\boldsymbol{H}}
\newcommand{\bmI}{\boldsymbol{I}}	\newcommand{\bmJ}{\boldsymbol{J}}
\newcommand{\bmK}{\boldsymbol{K}}	\newcommand{\bmL}{\boldsymbol{L}}
\newcommand{\bmM}{\boldsymbol{M}}	\newcommand{\bmN}{\boldsymbol{N}}
\newcommand{\bmO}{\boldsymbol{O}}	\newcommand{\bmP}{\boldsymbol{P}}
\newcommand{\bmQ}{\boldsymbol{Q}}	\newcommand{\bmR}{\boldsymbol{R}}
\newcommand{\bmS}{\boldsymbol{S}}	\newcommand{\bmT}{\boldsymbol{T}}
\newcommand{\bmU}{\boldsymbol{U}}	\newcommand{\bmV}{\boldsymbol{V}}
\newcommand{\bmW}{\boldsymbol{W}}	\newcommand{\bmX}{\boldsymbol{X}}
\newcommand{\bmY}{\boldsymbol{Y}}	\newcommand{\bmZ}{\boldsymbol{Z}}
%Small Letters
\newcommand{\bma}{\boldsymbol{a}}	\newcommand{\bmb}{\boldsymbol{b}}
\newcommand{\bmc}{\boldsymbol{c}}	\newcommand{\bmd}{\boldsymbol{d}}
\newcommand{\bme}{\boldsymbol{e}}	\newcommand{\bmf}{\boldsymbol{f}}
\newcommand{\bmg}{\boldsymbol{g}}	\newcommand{\bmh}{\boldsymbol{h}}
\newcommand{\bmi}{\boldsymbol{i}}	\newcommand{\bmj}{\boldsymbol{j}}
\newcommand{\bmk}{\boldsymbol{k}}	\newcommand{\bml}{\boldsymbol{l}}
\newcommand{\bmm}{\boldsymbol{m}}	\newcommand{\bmn}{\boldsymbol{n}}
\newcommand{\bmo}{\boldsymbol{o}}	\newcommand{\bmp}{\boldsymbol{p}}
\newcommand{\bmq}{\boldsymbol{q}}	\newcommand{\bmr}{\boldsymbol{r}}
\newcommand{\bms}{\boldsymbol{s}}	\newcommand{\bmt}{\boldsymbol{t}}
\newcommand{\bmu}{\boldsymbol{u}}	\newcommand{\bmv}{\boldsymbol{v}}
\newcommand{\bmw}{\boldsymbol{w}}	\newcommand{\bmx}{\boldsymbol{x}}
\newcommand{\bmy}{\boldsymbol{y}}	\newcommand{\bmz}{\boldsymbol{z}}

%---------------------------------------
% Scr Math Fonts :-
%---------------------------------------

\newcommand{\sA}{{\mathscr{A}}}   \newcommand{\sB}{{\mathscr{B}}}
\newcommand{\sC}{{\mathscr{C}}}   \newcommand{\sD}{{\mathscr{D}}}
\newcommand{\sE}{{\mathscr{E}}}   \newcommand{\sF}{{\mathscr{F}}}
\newcommand{\sG}{{\mathscr{G}}}   \newcommand{\sH}{{\mathscr{H}}}
\newcommand{\sI}{{\mathscr{I}}}   \newcommand{\sJ}{{\mathscr{J}}}
\newcommand{\sK}{{\mathscr{K}}}   \newcommand{\sL}{{\mathscr{L}}}
\newcommand{\sM}{{\mathscr{M}}}   \newcommand{\sN}{{\mathscr{N}}}
\newcommand{\sO}{{\mathscr{O}}}   \newcommand{\sP}{{\mathscr{P}}}
\newcommand{\sQ}{{\mathscr{Q}}}   \newcommand{\sR}{{\mathscr{R}}}
\newcommand{\sS}{{\mathscr{S}}}   \newcommand{\sT}{{\mathscr{T}}}
\newcommand{\sU}{{\mathscr{U}}}   \newcommand{\sV}{{\mathscr{V}}}
\newcommand{\sW}{{\mathscr{W}}}   \newcommand{\sX}{{\mathscr{X}}}
\newcommand{\sY}{{\mathscr{Y}}}   \newcommand{\sZ}{{\mathscr{Z}}}


%---------------------------------------
% Math Fraktur Font
%---------------------------------------

%Captital Letters
\newcommand{\mfA}{\mathfrak{A}}	\newcommand{\mfB}{\mathfrak{B}}
\newcommand{\mfC}{\mathfrak{C}}	\newcommand{\mfD}{\mathfrak{D}}
\newcommand{\mfE}{\mathfrak{E}}	\newcommand{\mfF}{\mathfrak{F}}
\newcommand{\mfG}{\mathfrak{G}}	\newcommand{\mfH}{\mathfrak{H}}
\newcommand{\mfI}{\mathfrak{I}}	\newcommand{\mfJ}{\mathfrak{J}}
\newcommand{\mfK}{\mathfrak{K}}	\newcommand{\mfL}{\mathfrak{L}}
\newcommand{\mfM}{\mathfrak{M}}	\newcommand{\mfN}{\mathfrak{N}}
\newcommand{\mfO}{\mathfrak{O}}	\newcommand{\mfP}{\mathfrak{P}}
\newcommand{\mfQ}{\mathfrak{Q}}	\newcommand{\mfR}{\mathfrak{R}}
\newcommand{\mfS}{\mathfrak{S}}	\newcommand{\mfT}{\mathfrak{T}}
\newcommand{\mfU}{\mathfrak{U}}	\newcommand{\mfV}{\mathfrak{V}}
\newcommand{\mfW}{\mathfrak{W}}	\newcommand{\mfX}{\mathfrak{X}}
\newcommand{\mfY}{\mathfrak{Y}}	\newcommand{\mfZ}{\mathfrak{Z}}
%Small Letters
\newcommand{\mfa}{\mathfrak{a}}	\newcommand{\mfb}{\mathfrak{b}}
\newcommand{\mfc}{\mathfrak{c}}	\newcommand{\mfd}{\mathfrak{d}}
\newcommand{\mfe}{\mathfrak{e}}	\newcommand{\mff}{\mathfrak{f}}
\newcommand{\mfg}{\mathfrak{g}}	\newcommand{\mfh}{\mathfrak{h}}
\newcommand{\mfi}{\mathfrak{i}}	\newcommand{\mfj}{\mathfrak{j}}
\newcommand{\mfk}{\mathfrak{k}}	\newcommand{\mfl}{\mathfrak{l}}
\newcommand{\mfm}{\mathfrak{m}}	\newcommand{\mfn}{\mathfrak{n}}
\newcommand{\mfo}{\mathfrak{o}}	\newcommand{\mfp}{\mathfrak{p}}
\newcommand{\mfq}{\mathfrak{q}}	\newcommand{\mfr}{\mathfrak{r}}
\newcommand{\mfs}{\mathfrak{s}}	\newcommand{\mft}{\mathfrak{t}}
\newcommand{\mfu}{\mathfrak{u}}	\newcommand{\mfv}{\mathfrak{v}}
\newcommand{\mfw}{\mathfrak{w}}	\newcommand{\mfx}{\mathfrak{x}}
\newcommand{\mfy}{\mathfrak{y}}	\newcommand{\mfz}{\mathfrak{z}}


\newcommand{\mytitle}{ECO\_3TC41\_TP Introduction à l'économie}
\newcommand{\myauthor}{github/omgeta}
\newcommand{\mydate}{AU 26-27 S1}

\begin{document}
\raggedright
\footnotesize
\begin{multicols*}{3}
\setlength{\premulticols}{1pt}
\setlength{\postmulticols}{1pt}
\setlength{\multicolsep}{1pt}
\setlength{\columnsep}{2pt}

{\normalsize{\textbf{\mytitle}}}\\
{\footnotesize{\mydate\hspace{2pt}\textemdash\hspace{2pt}\myauthor}}
%%%%%%%%%%%%%%%%%%%%%%%%%%%%%%%%%%%%%%%%%%%%%%%%%%%%%%
%                      Begin                         %
%%%%%%%%%%%%%%%%%%%%%%%%%%%%%%%%%%%%%%%%%%%%%%%%%%%%%%
\section{Introduction}

L'économie étudie les choix des agents économiques, l'allocation des ressources et le fonctionnement des marchés. Grandes questions:
\begin{enumerate}[\roman*.]
  \item IA et Transition Numérique: croissance, innovation, compétences, inégalités, protection de la vie privée 
  \item Climat: pollution, coordination internationale
  \item Inflation: pouvoir de marché et concurrence, bulles financières, pouvoir d’achat, impôts, politiques monétaires
\end{enumerate}

Circuit économique:
\begin{enumerate}[\roman*.]
  \item Ménages: consomment les biens et services, fournissent les facteurs de production
  \item Entreprises: produires les biens et services, utilisent les facteurs de production
\end{enumerate}

Complexités:
\begin{enumerate}[\roman*.]
  \item Tout est interconnecté: equilibre général (e.g. 2008)
  \item Futur détermine le présent: anticipations, arbitrage intertemporel
  \item Asymétries d'information
\end{enumerate}

\textbf{Ex. Café en gare/aéroport (prix élevé)}:
\begin{enumerate}[\roman*.]
  \item Marge (markup): prix $-$ coût marginal
  \item Monopole local: absence de concurrents, demande captive
  \item Choix endogènes: emplacement stratégique, contrats d'exclusivité
  \item Marché contestable: faibles coûts fixes d'entrée, mais accès limité aux emplacements
  \item Extraction du surplus: propriétaire des locaux capte les profits des kiosques via les loyers
\end{enumerate}
Conclusion: prix élevé principalement dû à la rareté des espaces commerciaux et au pouvoir du propriétaire, plutôt qu'à la seule demande.

\colbreak
\section{Transactions}

\subsection{Comment organiser les transactions ?}

Approches pour déterminer les prix:
\begin{enumerate}[\roman*.]
  \item Décentralisée: négociations bi/multilatérales (selon client et pouvoir de négociation) ou automatisées (énergie P2P, prêts P2P)
  \item Centralisée: planification (Marx) ou marché (concurrence parfaite)
\end{enumerate}

Négociations bilatérales:
\begin{itemize}[$-$]
  \item Coûts de transaction (Coase): rechercher, négocier et contrôler chaque échange est coûteux ; les entreprises internalisent les transactions si coûte moins cher, expliquant leur taille 
  \item Contrats incomplets: ne prévoient pas toutes les situations avec risques informationnels 
    \begin{itemize}[topsep=0.5em]
      \item Sélection adverse: information cachée avant
      \item Aléas moral: comportement caché après 
    \end{itemize}
  \item N'exploitent pas les économies d'échelle (coût moyen diminue quand production augmente)
\end{itemize}

\subsection{Marx vs Hayek}

Coordination économique:
\begin{enumerate}[\roman*.]
  \item Marx: centralisée ; une autorité organise production, allocation et distribution des ressources
  \item Hayek: décentralisée par le marché ; prix transmettent et agrègent l'information dispersée entre agents, sans planificateur central
\end{enumerate}

Marchés et limites de la planification:
\begin{enumerate}[\roman*.]
  \item École de Chicago: importance des marchés, des prix et des incitations individuelles
  \item Valeur attribuée à un bien $\neq$ Prix d'échange
  \item IA et planification: même avec une forte capacité de calcul, problèmes d'informations privées et de critères de redistribution
\end{enumerate}

\colbreak
\section{Marché Concurrentiel}

\subsection{La demande et l'offre}

\textbf{Demande individuelle}: quantité qu'un consommateur souhaite et peut acheter à chaque prix.

\textbf{Courbes d'indifférence}: ensembles de paniers $(x,y)$ procurant la même utilité $U(x,y)$.
\begin{enumerate}[\roman*.]
  \item Plus éloignées de l'origine $\Rightarrow$ utilité supérieure
  \item Ne se croisent pas
  \item Décroissantes et convexes (TMS décroissant)
\end{enumerate}

\textbf{Taux Marginal de Substitution (TMS)}: quantité de $y$ sacrifiée pour une unité supplémentaire de $x$ à utilité constante.
\[
TMS_{x,y}=-\left.\frac{dy}{dx}\right|_{U=\bar U}
=\frac{U'_x}{U'_y}
\]
\begin{enumerate}[\roman*.]
  \item Décroissant: plus on possède de $x$, moins on accepte de sacrifier de $y$ pour en obtenir davantage.
\end{enumerate}

\textbf{Optimum du consommateur}: maximiser l'utilité sous contrainte budgétaire:
\[
P_xx+P_yy\leq R,\qquad
\boxed{TMS_{x,y}=\frac{P_x}{P_y}}
\]
\begin{enumerate}[\roman*.]
  \item À l'optimum intérieur: tangence entre courbe d'indifférence et droite budgétaire (pente $-P_x/P_y$).
\end{enumerate}

\textbf{Effets sur la consommation}:
\begin{enumerate}[\roman*.]
  \item Effet de revenu: variation du pouvoir d'achat ; bien normal ($R\uparrow\Rightarrow Q_d\uparrow$), bien inférieur ($R\uparrow\Rightarrow Q_d\downarrow$)
  \item Effet de substitution: variation des prix relatifs ; consommation du bien devenu moins cher $\uparrow$
\end{enumerate}

\textbf{Courbe de demande}: obtenue à partir des paniers optimaux pour différents prix ; généralement décroissante.

\textbf{Demande agrégée}: somme horizontale des demandes individuelles au même prix:
\[
Q_D(P)=\sum_{i=1}^{n}Q_{d,i}(P)
\]

\colbreak

\textbf{Offre individuelle}: quantité qu'un producteur souhaite vendre à chaque prix. 

\textbf{Courbe d'offre}: généralement croissante ; prix plus élevé incite à produire davantage.

\textbf{Offre agrégée}: somme horizontale des offres individuelles au même prix:
\[
Q_S(P)=\sum_{i=1}^{n}Q_{s,i}(P)
\]

\textbf{Élasticité-prix}: sensibilité relative des quantités demandées/offertes à une variation du prix.
\[
\boxed{\varepsilon_d=-\frac{\%\Delta Q_d}{\%\Delta P}
=-\frac{P}{Q_d}\frac{dQ_d}{dP},\quad \varepsilon_s=\frac{\%\Delta Q_s}{\%\Delta P}
=\frac{P}{Q_s}\frac{dQ_s}{dP}}
\]
\begin{enumerate}[\roman*.]
  \item $\varepsilon=0$: parfaitement inélastique (verticale)
  \item $0<\varepsilon<1$: inélastique
  \item $\varepsilon=1$: unitaire
  \item $\varepsilon>1$: élastique
  \item $\varepsilon\to\infty$: parfaitement élastique (horizontale)
\end{enumerate}

\textbf{Déplacement de la demande}:
\begin{enumerate}[\roman*.]
  \item Variation du prix propre: mouvement le long de la courbe
  \item Autres facteurs: revenu, goûts, anticipations, nombre d'acheteurs, prix des biens liés
  \item Biens substituts: $P_y\uparrow\Rightarrow D_x\uparrow$
  \item Biens complémentaires: $P_y\uparrow\Rightarrow D_x\downarrow$
\end{enumerate}
Demande $\uparrow$: courbe vers la droite ; demande $\downarrow$: vers la gauche.

\textbf{Déplacement de l'offre}:
\begin{enumerate}[\roman*.]
  \item Variation du prix propre: mouvement le long de la courbe
  \item Coûts des inputs $\uparrow\Rightarrow$ offre $\downarrow$
  \item Progrès technologique $\Rightarrow$ offre $\uparrow$
  \item Anticipations des prix futurs
  \item Nombre de vendeurs $\uparrow\Rightarrow$ offre de marché $\uparrow$
\end{enumerate}
Offre $\uparrow$: courbe vers la droite ; offre $\downarrow$: vers la gauche.

\subsection{L'équilibre concurrentiel}

\textbf{Équilibre de marché}: intersection de l'offre et de la demande à prix d'équilibre $P^*$:
\[
\boxed{Q_D(P^*)=Q_S(P^*)=Q^*}
\]
\begin{enumerate}[\roman*.]
  \item Excédent: $Q_S>Q_D$ $\Rightarrow$ pression à la baisse sur $P$
  \item Pénurie: $Q_D>Q_S$ $\Rightarrow$ pression à la hausse sur $P$
\end{enumerate}

\textbf{Déplacements de l'équilibre}:
\begin{enumerate}[\roman*.]
  \item Demande $\uparrow$: $P^*\uparrow$, $Q^*\uparrow$
  \item Demande $\downarrow$: $P^*\downarrow$, $Q^*\downarrow$
  \item Offre $\uparrow$: $P^*\downarrow$, $Q^*\uparrow$
  \item Offre $\downarrow$: $P^*\uparrow$, $Q^*\downarrow$
\end{enumerate}

\textbf{Contrôle des prix}:
\begin{enumerate}[\roman*.]
  \item Prix plafond: prix maximal légal
    \begin{itemize}[topsep=0.5em]
      \item si $P_{\max}<P^*$ $\Rightarrow$ pénurie ($Q_D>Q_S$) 
    \end{itemize}
  \item Prix plancher: prix minimal légal
    \begin{itemize}[topsep=.5em]
      \item si $P_{\min}>P^*$ $\Rightarrow$ excédent ($Q_S>Q_D$)
    \end{itemize}
  \item Contrôle des loyers: pénurie potentiellement accrue à long terme, lorsque offre et demande deviennent plus sensibles au prix
  \item Salaire minimum: dans le modèle concurrentiel, salaire au-dessus de l'équilibre $\Rightarrow$ excès d'offre de travail (chômage)
  \item Rationnement: répartition non tarifaire des biens en situation de pénurie
\end{enumerate}

\textbf{Taxes unitaires}: taxe fixe $t$ sur chaque unité vendue ; augmente le prix payé par l'acheteur et diminue le prix reçu par le vendeur:
\[
\boxed{P_a-P_v=t}
\]
\begin{enumerate}[\roman*.]
  \item Taxe sur vendeurs: augmente le coût de vente $\Rightarrow$ offre se décale vers le haut
  \item Taxe sur acheteurs: réduit leur disposition à payer hors taxe $\Rightarrow$ demande se décale vers le bas
  \item Dans les deux cas: $P_a\uparrow$, $P_v\downarrow$, $Q^*\downarrow$ ; même résultat quel que soit celui qui verse la taxe à l'État
\end{enumerate}

\textbf{Incidence fiscale}: partage réel de la taxe entre acheteurs et vendeurs, selon leur sensibilité au prix (élasticité).
\begin{enumerate}[\roman*.]
  \item Demande peu élastique: acheteurs continuent d'acheter malgré hausse du prix $\Rightarrow$ supportent davantage la taxe
  \item Offre peu élastique: vendeurs réduisent peu leur production malgré baisse du prix reçu $\Rightarrow$ supportent davantage la taxe
\end{enumerate}
Règle: côté le moins élastique supporte la plus grande part de la taxe.

\textbf{Limites de la fiscalité}:
\begin{enumerate}[\roman*.]
  \item Nozick (Wilt Chamberlain): même après redistribution égalitaire, échanges volontaires peuvent recréer des inégalités ; question du droit de taxer pour redistribuer
  \item Vote par les pieds: individus peuvent déménager vers des territoires aux impôts/services publics plus avantageux
  \item Courbe de Laffer: recettes $=$ taux d'impôt $\times$ base imposable ; si taux trop élevé, activité imposable peut diminuer suffisamment pour réduire les recettes
\end{enumerate}

\colbreak
\section{Bien-être}

\subsection{Surplus du consommateur}

\textbf{Disponibilité à payer (DAP)}: prix maximal qu'un acheteur accepte de payer pour un bien ; représente la valeur qu'il lui attribue.

\textbf{Surplus du consommateur (SC)}: gain obtenu en payant moins que sa DAP.
\[
\boxed{SC=DAP-P}
\]
\textbf{Courbe de demande}: sa hauteur représente la DAP marginale. Le SC total est l'aire sous la demande et au-dessus du prix, jusqu'à la quantité achetée.
\[
SC=\frac12(P_{\max}-P)Q
\]
(pour une demande linéaire)

Prix $\downarrow\Rightarrow SC\uparrow$: acheteurs existants paient moins et de nouveaux acheteurs entrent sur le marché.

\subsection{Surplus du producteur}

\textbf{Disponibilité à vendre}: prix minimal qu'un producteur accepte, déterminé par son coût de production.

\textbf{Surplus du producteur (SP)}: gain obtenu en vendant à un prix supérieur au coût.
\[
\boxed{SP=P-C}
\]
\textbf{Courbe d'offre}: sa hauteur représente le coût marginal. Le SP total est l'aire au-dessus de l'offre et sous le prix, jusqu'à la quantité vendue.
\[
SP=\frac12(P-P_{\min})Q
\]
(pour une offre linéaire)

Prix $\uparrow\Rightarrow SP\uparrow$: producteurs existants gagnent davantage et de nouveaux producteurs peuvent entrer.

\subsection{Surplus total}

\textbf{Surplus total (ST)}: somme des gains des consommateurs et des producteurs:
\[
\boxed{ST=SC+SP=DAP_{\text{totale}}-C_{\text{total}}}
\]

\textbf{Efficacité du marché}: l'équilibre concurrentiel maximise le surplus total (en absence de défaillances de marché).
\begin{enumerate}[\roman*.]
  \item $Q<Q^*$: DAP marginale $>$ coût marginal $\Rightarrow$ échange supplémentaire augmente le surplus
  \item $Q>Q^*$: coût marginal $>$ DAP marginale $\Rightarrow$ échange supplémentaire réduit le surplus
  \item $Q=Q^*$: DAP marginale $=$ coût marginal $\Rightarrow ST$ maximal
\end{enumerate}

\incimg{surplus}
\colbreak
\subsection{Efficacité de Pareto}

\textbf{Amélioration de Pareto}: améliore le bien-être d'au moins un agent sans détériorer celui des autres.
\textbf{Optimum de Pareto}: allocation où aucune amélioration de Pareto supplémentaire n'est possible

\textbf{Boîte d'Edgeworth}: répartition de deux biens entre deux agents. À l'optimum intérieur, leurs courbes d'indifférence sont tangentes:
\[
\boxed{TMS_A=TMS_B}
\]
\begin{center}\vspace{-1em}
  \incimg[0.7]{edgeworth}
\end{center}\vspace{-1em}
\textbf{Courbe des contrats}: ensemble des allocations Pareto optimales ; échanges depuis l'allocation initiale vers une allocation avantageuse pour les deux agents.

\textbf{Product Mix}: combinaison efficace des biens produits ; le TMT mesure la quantité d'un bien sacrifiée pour produire davantage de l'autre.
\[
\boxed{TMS=TMT}
\]
Préférences des consommateurs = possibilités de production.

\textbf{Théorèmes du bien-être}:
\begin{enumerate}[\arabic*.]
  \item équilibre concurrentiel $\Rightarrow$ optimum de Pareto (main invisible)
  \item optimum de Pareto $\Rightarrow$ équilibre concurrentiel possible après redistribution initiale
\end{enumerate}
Sous certaines hypothèses: concurrence, absence de défaillances de marché. Redistribution par impôt forfaitaire pour éviter de modifier les incitations.

\textbf{Efficacité statique}: allocation optimal à instant donné ; \textbf{efficacité dynamique}: prise en compte de l'innovation et des investissements futurs.
\vspace{-1em}
\colbreak
\subsection{Critiques du critère de Pareto}

\textbf{Limite principale}: efficacité $\neq$ équité ; une allocation très inégalitaire peut être Pareto optimale.

\textbf{Amartya Sen}: bien-être ne se limite pas à l'utilité/consommation ; inclut libertés individuelles, droits politiques, environnement et participation sociale.

\subsection{L'approche de John Rawls}

\textbf{Voile de l'ignorance}: choisir les règles d'une société sans connaître sa future position sociale, richesse, religion ou génération.

\textbf{Principes de justice}:
\begin{enumerate}[\roman*.]
  \item Libertés fondamentales égales pour tous
  \item Égalité équitable des chances
  \item Principe de différence: inégalités acceptables si elles bénéficient aux plus défavorisés
\end{enumerate}

\textbf{Critère max-min}: privilégier l'amélioration de la situation des individus les moins favorisés.

\textbf{Critiques}: redistribution fiscale peut réduire les incitations au travail/investissement/innovation ; justice distributive et efficacité économique pas nécessairement compatibles.

\textbf{Équité algorithmique}: critères de justice indépendants de l'identité des personnes ; robustesse aux permutations des individus.

\colbreak
\section{Long terme}

\textbf{Court terme}: certains facteurs de production sont fixes ; \textbf{long terme}: tous les facteurs peuvent être ajustés, entreprises peuvent entrer/sortir du marché.

\textbf{Coûts économiques}:
\begin{enumerate}[\roman*.]
  \item Coûts explicites: dépenses directes\\(ex. salaires, loyers, matières premières)
  \item Coûts implicites: valeur des ressources utilisées sans paiement direct
  \item Coût d'opportunité: valeur de la meilleure alternative abandonnée\\(ex. études: salaire non perçu)
\end{enumerate}

\textbf{Profits}:
\begin{enumerate}[\roman*.]
  \item Comptable: recettes $-$ coûts explicites
  \item Économique: recettes $-$ coûts explicites $-$ coûts implicites
\end{enumerate}

\textbf{Coût total ($CT$)} = Coût fixe + Coût variable 
\begin{enumerate}[\roman*.]
  \item Coût fixe ($CF$): indépendant de la quantité
  \item Coût variable ($CV$): varie avec la production
\end{enumerate}

\textbf{Coûts moyens}: coût de production par unité
\[
\boxed{
CFM=\frac{CF}{Q},\quad
CVM=\frac{CV}{Q},\quad
CTM=\frac{CT}{Q}}
\]

\textbf{Coût marginal ($Cm$)}: coût d'une unité supplémentaire
\[
\boxed{Cm=\frac{\Delta CT}{\Delta Q}}
\]

\textbf{Relations entre les courbes}:
\begin{enumerate}[\roman*.]
  \item $CFM$ décroissant (coûts fixes répartis sur davantage d'unités)
  \item $CVM$ et $CTM$ généralement en U
  \item $Cm<CTM\Rightarrow CTM\downarrow$
  \item $Cm>CTM\Rightarrow CTM\uparrow$
  \item $Cm=CTM\Rightarrow CTM$ minimal (courbe en U)
\end{enumerate}

\textbf{Rendements d'échelle}: variation de production lorsque tous les facteurs augmentent dans la même proportion.
\begin{enumerate}[\roman*.]
  \item Croissants: facteurs $\times2\Rightarrow$ production $> \times2$
  \item Constants: facteurs $\times2\Rightarrow$ production $\times2$
  \item Décroissants: facteurs $\times2\Rightarrow$ production $< \times2$
\end{enumerate}

\textbf{Économies d'échelle}: coût moyen de long terme diminue quand la production augmente ; \textbf{déséconomies}: coût moyen augmente (ex. difficultés de coordination).

{\centering\incimg[0.7]{rendements}\par}
\vspace{-1em}
\subsection{Équilibre de long terme}

\textbf{Concurrence parfaite}: entreprise preneuse de prix ; maximise son profit en produisant jusqu'à $P=Cm$.

\textbf{Entrée et sortie}:
\begin{enumerate}[\roman*.]
  \item $P>CTM$: profit économique positif \\$\Rightarrow$ entrée de concurrents, offre $\uparrow$, prix $\downarrow$
  \item $P<CTM$: pertes \\$\Rightarrow$ sortie d'entreprises, offre $\downarrow$, prix $\uparrow$
  \item $P=CTM$: profit économique nul \\$\Rightarrow$ équilibre de long terme
\end{enumerate}
{\centering\incimg[0.7]{equilibre}\par}
\vspace{-1em}
\colbreak
\section{Défaillances}

\textbf{Défaillance de marché}: situation où le marché n'alloue pas efficacement les ressources, entraînant une perte de bien-être collectif.

\subsection{Monopole}

\textbf{Monopole}: un seul vendeur, sans substitut proche ; dispose d'un pouvoir sur le prix grâce aux barrières à l'entrée:
\begin{enumerate}[\roman*.]
  \item Monopolisation d'un input clé 
  \item Régulation accordant une exclusivité
  \item Monopole naturel: économies d'échelle
  \item Propriété intellectuelle
\end{enumerate}

\textbf{Maximisation du profit}: monopole choisit la quantité maximisant recettes moins coûts:
\[
\Pi(Q)=RT(Q)-CT(Q),\qquad \boxed{Rm=Cm}
\]

\begin{center}
  \incimg[0.8]{monopole}
\end{center}

\textbf{Perte sèche}: surplus collectif perdu car le monopole produit moins et vend plus cher qu'en concurrence.
\begin{enumerate}[\roman*.]
  \item Monopole: $Rm=Cm$, quantité $Q_M$
  \item Optimum social: $P=Cm$, quantité $Q^*$
  \item $Q_M<Q^*$ et $P_M>P^*$ $\Rightarrow$ échanges mutuellement avantageux non réalisés
\end{enumerate}

\colbreak

\subsection{Externalités}

\textbf{Externalité}: effet d'une activité sur des tiers, non pris en compte dans les prix du marché

\textbf{Externalité négative} (ex. aluminium): producteurs ignorent le coût de la pollution imposé à la société.
\[
\boxed{CMS=CMP+CME}
\]
\begin{enumerate}[\roman*.]
  \item Coût marginal social = privé + externe ;\\ $Q_{\text{marché}}>Q_{\text{optimal}}$ $\Rightarrow$ surproduction.
\end{enumerate}

\textbf{Externalité positive} (ex. éducation): consommateurs ignorent une partie des bénéfices apportés à la société.
\[
\boxed{BMS=BMP+BME}
\]
\begin{enumerate}[\roman*.]
  \item Bénéfice marginal social = privé + externe ;\\ $Q_{\text{marché}}<Q_{\text{optimal}}$ $\Rightarrow$ sous-consommation.
\end{enumerate}

\textbf{Solutions aux externalités}:
\begin{enumerate}[\roman*.]
  \item Coase: négociation entre agents si droits de propriété définis et coûts de transaction faibles ; internalisation de l'externalité
  \item Réglementation: interdictions/limites directes
  \item Taxe pigouvienne: pollueur paie le dommage 
  \item Cap and trade: plafond total d'émissions et quotas échangeables (ex. EU ETS)
\end{enumerate}

\textbf{Limites de la taxe pigouvienne}:
\begin{enumerate}[\roman*.]
  \item Effets redistributifs: ménages modestes davantage affectés (ex. Gilets jaunes)
  \item Pollution et dommages difficiles à mesurer
  \item Difficulté d'estimer la disponibilité à payer pour l'environnement
\end{enumerate}

\textbf{Taxe vs quotas}:
\begin{enumerate}[\roman*.]
  \item Taxe: fixe le prix de la pollution ; quantité s'ajuste
  \item Quotas: fixe la quantité de pollution ; prix s'ajuste
\end{enumerate}

\textbf{Limites des quotas}: droit de polluer légitimé, fraude, manipulation/collusion, délocalisation vers pays moins réglementés (forum shopping), incitations insuffisantes à l'innovation.

\subsection{Les biens publics}

\textbf{Typologie des biens}:
\begin{enumerate}[\roman*.]
  \item Bien privé: rival, excluable (ex. glace)
  \item Bien de club: non rival, excluable\\(ex. autoroute non congestionnée à péage)
  \item Bien commun: rival, non excluable (ex. poissons)
  \item Bien public pur: non rival, non excluable\\(ex. défense nationale)
\end{enumerate}

\textbf{Biens publics impurs}: variables selon l'utilisation.
\begin{enumerate}[\roman*.]
  \item Congestion: bien devient rival si trop utilisé\\(ex. hôpitaux pendant COVID)
  \item Coût d'exclusion: coût pour empêcher l'accès aux non-payeurs (ex. péage)
\end{enumerate}

\textbf{Passager clandestin}: agent profite d'un bien sans contribuer à son financement ; entreprises ne récupèrent pas tous les bénéfices $\Rightarrow$ sous-investissement privé\\(ex. recherche fondamentale).

\subsection{Les biens communs}

\textbf{Tragédie des communs}: ressource collective surexploitée car chaque agent privilégie son gain individuel sans supporter tous les coûts (ex. surpêche, surpâturage). Solutions:
\begin{enumerate}[\roman*.]
  \item Ostrom: règles collectives, surveillance
  \item Tragédie des anti-communs: trop de droits d'exclusion empêchent l'utilisation efficace d'une ressource $\Rightarrow$ sous-exploitation 
\end{enumerate}

\subsection{Les problèmes informationnels}

\textbf{Ex. Akerlof (lemons)}: 50\% plums, 50\% lemons
\begin{enumerate}[\roman*.]
  \item Vendeurs: minimum \$2000 (plum), \$1000 (lemon)
  \item Acheteurs: DAP \$2400 (plum), \$1200 (lemon)
  \item Qualité inconnue $\Rightarrow$ DAP moyenne $=0.5(2400)+0.5(1200)=\$1800$
  \item Vendeurs de plums exigent \$2000 $\Rightarrow$ quittent
  \item Restent uniquement les lemons $\Rightarrow$ DAP tombe à \$1200 ; disparition du marché des bonnes voitures
\end{enumerate}
\vspace{-1em}
\colbreak
\section{Rationalité}

\textbf{Rationalité économique}: choix cohérents avec les préférences de l'agent ; représentés par fonction $u(x)$

\textbf{Relations de préférence}:
\begin{enumerate}[\roman*.]
  \item $x\succ y$: $x$ strictement préféré à $y$
  \item $x\sim y$: indifférence entre $x$ et $y$
  \item $x\succeq y$: $x$ au moins aussi bon que $y$
\end{enumerate}

\textbf{Axiomes de rationalité}:
\begin{enumerate}[\roman*.]
  \item Complétude: tous comparables $x\succeq y$ ou $y\succeq x$
  \item Réflexivité: 
  $x\succeq x$
  \item Transitivité: 
  $x\succeq y,\ y\succeq z\Rightarrow x\succeq z$
\end{enumerate}

\textbf{Indépendance of Irrelevant Alternatives (IIA)}: ajouter une option ne devrait pas modifier la préférence entre deux options existantes

\textbf{Biais cognitifs}: décisions influencées par le contexte, incohérentes avec rationalité théorique (ex. dark patterns)


\subsection{Préférences révélées}

\textbf{Préférences révélées}: déduire les préférences à partir des choix observés, plutôt que des déclarations.

\textbf{Contrainte budgétaire}:
\[
p_1x_1+p_2x_2\leq m
\]
Si $x$ est choisi alors que $y$ est abordable, $x$ est révélé préféré à $y$.

\textbf{WARP (axiome faible)}: cohérence des choix entre différents budgets.
\begin{enumerate}[\roman*.]
  \item Vérifié: $x$ choisi quand $y$ abordable ; si $y$ choisi ensuite, $x$ n'est plus abordable
  \item Violé: $x$ choisi quand $y$ abordable, puis $y$ choisi quand $x$ est encore abordable (choix distincts)
\end{enumerate}
WARP permet de tester la rationalité à partir des comportements observés.

\colbreak
\subsection{Théorie des jeux}

\textbf{Jeu stratégique}: situation où le gain d'un joueur dépend de ses décisions et de celles des autres.
\begin{enumerate}[\roman*.]
  \item Joueurs: agents prenant les décisions
  \item Stratégies: actions disponibles
  \item Gains (payoffs): utilités des résultats
  \item Actions successives et simultanées 
\end{enumerate}

\textbf{Dilemme du prisonnier}: chaque joueur a intérêt à trahir, quelle que soit l'action de l'autre ; les deux trahissent, alors que coopérer ensemble leur serait plus avantageux. Rationalité individuelle $\neq$ optimum collectif.

\textbf{Équilibre de Nash}: aucun joueur ne peut améliorer son gain en changeant seul de stratégie, connaissant celles des autres.
\[
\boxed{u_i(a_i^*,a_{-i}^*)\geq
u_i(a_i,a_{-i}^*)}
\]
Pour tout joueur $i$ et toute stratégie alternative $a_i$.

\textbf{Connaissance commune}: chacun sait que les autres savent une information, et ainsi de suite.

\textbf{Bach ou Stravinsky}: jeu de coordination ; les joueurs préfèrent être ensemble mais pas le même concert $\Rightarrow$ plusieurs équilibres de Nash possibles.

\textbf{Stratégies}:
\begin{enumerate}[\roman*.]
  \item Dominante: meilleure quelle que soit la décision des autres
  \item Dominée: toujours moins avantageuse qu'une autre stratégie
\end{enumerate}

\textbf{Trouver un équilibre de Nash}:
\begin{enumerate}[\roman*.]
  \item Identifier les meilleures réponses de chaque joueur aux stratégies adverses
  \item Éliminer les stratégies strictement dominées si possible
  \item Équilibre: chaque joueur joue une meilleure réponse aux autres
\end{enumerate}

\colbreak
\colbreak
\section{Concepts macroéconomiques}

\subsection{l'État et ses aggregats}

\textbf{État-gendarme}: limité aux fonctions régaliennes (loi, justice, police, défense, monnaie), financées par les impôts.\\
\textbf{État-providence}: intervention économique et sociale au-delà du régalien (ex. éducation, santé).

Fonctions socio-économiques:
\begin{enumerate}[\roman*.]
  \item Allocation: corriger les défaillances du marché par réglementation, services publics, concurrence
  \item Redistribution: réduire les inégalités via impôts, cotisations et prestations sociales
  \item Stabilisation: réduire les déséquilibres économiques par politiques budgétaire et monétaire
\end{enumerate}

\textbf{Produit Intérieur Brut (PIB)}: valeur des finaux produits par les unités résidentes pendant une période: 
\begin{enumerate}[\roman*.]
  \item Production: somme des valeurs ajoutées (VA)
    \begin{itemize}[topsep=.5em, leftmargin=*]
      \item VA: production $-$ consommations intermédiaires 
    \end{itemize}
  \item Revenus: rémunérations du travail + excédent d'exploitation + impôts nets sur la production
  \item Demande: consommation + investissement + dépenses publiques + exportations nettes
\end{enumerate}
\[
\boxed{PIB=C+I+EX-IM}
\]
\textbf{Revenue National Brut (RNB)}: revenus primaires reçus par les résidents, y compris depuis l'étranger:
\[
\boxed{RNB=PIB+R_{\text{reçus RDM}}-R_{\text{versés RDM}}}
\]
\textbf{Enchaînement des agrégats}:
\begin{enumerate}[\roman*.]
  \item Revenu national disponible brut\\$=$ RNB $+$ transferts de richesse $-$ impôts 
  \item Épargne brute\\$=$ Revenu disponible $-$ consommation finale
  \item Besoin/capacité de financement national\\$=$ Eparge brute $-$ investissement 
\end{enumerate}

\colbreak
\subsection{Finances publiques}

\textbf{Budget de l'État}: recettes et dépenses des administrations centrales\\
\begin{itemize}[leftmargin=*]
  \item Solde budgétaire = Recettes $-$ Dépenses de l'État
\end{itemize}
\textbf{Budget public}: recettes des administrations APU (ex. état, sécurité sociale, regions) 
\begin{itemize}[leftmargin=*]
  \item Solde public = Recettes $-$ Dépenses publiques totales
\end{itemize}

\textbf{Recettes publiques}:
\begin{enumerate}[\roman*.]
  \item Fiscales: impôts sur le revenu, sociétés, TVA, TICPE, etc.
  \item Non fiscales: revenus du patrimoine public, amendes, privatisations
\end{enumerate}

\textbf{Prélèvements obligatoires (PO)}: tout versement par les particuliers/entreprises aux APU.
\begin{enumerate}[\roman*.]
  \item Directs: proportionnels au revenu\\(ex. le revenu des personnes, sociétés, les fortune) 
  \item Indirects: indépendants du revenu\\(ex. TVA, TIPP, taxe carbone) 
\end{enumerate}
\[
\boxed{Taux_{PO}=\frac{PO}{PIB}\times100}
\]
\textbf{Dépenses publiques}: fonctionnement (salaires, services), investissement (infrastructures), prestations sociales, subventions, défense, intérêts de la dette.

\textbf{Déficit vs dette}:
\begin{enumerate}[\roman*.]
  \item Déficit: dépenses $>$ recettes pendant une période 
  \item Dette publique: des emprunts non remboursés
  \item Déficits successifs $\Rightarrow$ dette généralement $\uparrow$ ; ratio dette/PIB peut néanmoins diminuer si PIB croît assez vite
\end{enumerate}

\textbf{Critères de Maastricht}: déficit public de référence $<3\%$ du PIB ; dette publique de référence $<60\%$ du PIB (avec exceptions et trajectoires d'ajustement).

\textbf{Financement du déficit}: emprunt, utilisation d'actifs, augmentation des recettes ou réduction des dépenses.

\subsection{Indicateurs économiques}

\textbf{Inflation}: hausse généralisée et durable des prix, mesurée par l'IPC/IPCH.
\[
\boxed{\pi_t=\frac{IPC_t-IPC_{t-1}}{IPC_{t-1}}\times100}
\]
\begin{enumerate}[\roman*.]
  \item Inflation: prix $\uparrow$ ; pouvoir d'achat de la monnaie $\downarrow$
  \item Désinflation: inflation ralentit mais prix continuent d'augmenter
  \item Déflation: baisse générale des prix ($\pi<0$)
\end{enumerate}

\textbf{Croissance économique}: augmentation du PIB réel
\[
\boxed{g_t=\frac{PIB_t-PIB_{t-1}}{PIB_{t-1}}\times100}
\]
\begin{enumerate}[\roman*.]
  \item PIB nominal: mesuré aux prix courants 
  \item PIB réel: mesuré à prix constants ; élimine l'effet des variations de prix
  \item Croissance réelle $\approx$ croissance nominale $-$ inflation
\end{enumerate}
\[
\boxed{1+g_{\text{réel}}=\frac{1+g_{\text{nominal}}}{1+\pi}}
\]

\textbf{Récession}: recul de l'activité économique ; souvent identifiée par 2 trimestres consécutifs de baisse du PIB réel

\colbreak

\textbf{Chômage (BIT)}: personne sans emploi, disponible sous deux semaines et recherchant activement un travail (ou emploi commençant prochainement).

\[
\boxed{Taux_{\text{chômage}}=
\frac{\text{Chômeurs BIT}}{\text{Population active}}\times100}
\]

\textbf{Population active}: personnes en emploi + chômeurs.

\textbf{Indicateurs du marché du travail}:
\[
\boxed{
\begin{aligned}
T_{\text{activité}}&=\frac{\text{Actifs}}{\text{Pop. en âge}}\\
T_{\text{emploi}}&=\frac{\text{Actifs occupés}}{\text{Pop. en âge}}
\end{aligned}}
\]
\begin{enumerate}[\roman*.]
  \item Conjoncturel: dû au ralentissement économique
  \item Structurel: inadéquation compétences/emplois
  \item Frictionnel: temps nécessaire pour trouver un nouvel emploi
  \item Halo du chômage: personnes souhaitant travailler mais ne satisfaisant pas tous les critères BIT
\end{enumerate}
\textbf{Plein emploi}: chômage non nul car chômage frictionnel (et éventuellement structurel) persiste.

\textbf{Balance commerciale (B\&S)}:
\[
\boxed{BC=EX-IM}
\]
\begin{enumerate}[\roman*.]
  \item Balance courante: inclut aussi revenus primaires et transferts courants avec le RDM
\end{enumerate}

\textbf{Taux de change}: valeur d'une monnaie exprimée dans une autre (ex. 1.25\$/€).
\begin{enumerate}[\roman*.]
  \item Appréciation de l'euro: € vaut plus de \$ $\Rightarrow$\\ exports désfavorisées, imports moins chères
  \item Dépréciation de l'euro: € vaut moins de \$ $\Rightarrow$\\ exports favorisées, imports plus chères
\end{enumerate}

\colbreak
\subsection{Stabilisateurs automatiques}

\textbf{Stabilisateurs automatiques}: mécanismes fiscaux et sociaux réduisant les fluctuations économiques sans décision politique nouvelle (contracyclique).

\begin{enumerate}[\roman*.]
  \item Expansion: recettes $\uparrow$ $\Rightarrow$ dépenses sociales $\downarrow$ $\Rightarrow$ solde public s'améliore $\Rightarrow$ demande modérée 
  \item Récession: recettes $\downarrow$ $\Rightarrow$ dépenses sociales $\uparrow$ $\Rightarrow$ solde public se dégrade $\Rightarrow$ demande soutenue 
\end{enumerate}

\subsection{Politiques conjoncturelles}

\textbf{Politiques discrétionnaires}: interventions volontaires de l'État pour réguler la conjoncture: 
\begin{enumerate}[\roman*.]
  \item Expansionniste (relance): stimuler activité et emploi en période de récession
  \item Restrictive (rigueur): freiner demande et inflation ou réduire les déficits
\end{enumerate}

\textbf{Politique budgétaire}: modification des dépenses (budget public) et des recettes (impôts):
\begin{enumerate}[\roman*.]
  \item Politique de demande (Keynes): dépenses public $\uparrow$ ou impôts ménages $\downarrow$ $\Rightarrow$ consommation $\uparrow$ $\Rightarrow$ production et emploi $\uparrow$
  \item Politique d'offre: impôts/coûts des entreprises $\downarrow$, subventions $\uparrow$ $\Rightarrow$ investissement et production $\uparrow$ 
\end{enumerate}

\textbf{Demande effective (Keynes)}: demande anticipée par les entreprises ; si débouchés prévus $\uparrow$, production, investissement et emploi peuvent augmenter.

\textbf{Multiplicateur keynésien}: hausse des dépenses publiques engendre une valeur de production supérieure 
\[
\boxed{\Delta Y=k\Delta G,\qquad k=\frac{1}{1-c}}
\]
$c$: propension marginale à consommer 

\textbf{Effet d'éviction}: dépenses publiques financées par emprunt $\uparrow$ $\Rightarrow$ demande de financement $\uparrow$ $\Rightarrow$ taux d'intérêt $\uparrow$ $\Rightarrow$ investissement privé $\downarrow$ ; réduit l'effet multiplicateur.

\colbreak
\textbf{Politique monétaire (BCE)}: assurer la stabilité des prix (inflation IPCH proche de 2\%/an) en contrôlant la masse monétaire via les taux d'intérêt.

\textbf{Politique conventionnelle}:
\begin{enumerate}[\roman*.]
  \item Expansionniste: taux directeurs $\downarrow$ $\Rightarrow$ offre de monnaie et crédits $\uparrow$ $\Rightarrow$ consommation/investissement $\uparrow$ $\Rightarrow$ croissance $\uparrow$
  \item Restrictive: taux directeurs $\uparrow$ $\Rightarrow$ crédits $\downarrow$ $\Rightarrow$ demande $\downarrow$ $\Rightarrow$ inflation freinée
\end{enumerate}

\textbf{3 canaux d'intervention de la BCE}:
\begin{enumerate}[\roman*.]
  \item Marché monétaire/interbancaire: BCE $\leftrightarrow$ banques commerciales ; modulation des 3 taux directeurs (refinancement, dépôt, prêt marginal) pour influencer les crédits
  \item Banque centrale $\leftrightarrow$ Trésor public: marché des titres de dette publique (bons du Trésor)
  \item Marché des changes: BCE $\leftrightarrow$ reste du monde ; interventions visant la stabilité du taux de change de l'euro
\end{enumerate}

\textbf{Policy mix}: combinaison des politiques budgétaire et monétaire ; une politique monétaire accommodante peut renforcer une relance budgétaire en limitant l'effet d'éviction.



\colbreak
\section{Politiques conjoncturelles}

\subsection{Relance (1970--1982)}

\textbf{Contexte international}:
\begin{enumerate}[\roman*.]
  \item Fin de Bretton Woods (1971): fin de la convertibilité du dollar en or $\Rightarrow$ changes flexibles dès 1973
  \item Serpent monétaire européen (1972--1979): fluctuations limitées à $\pm2,25\%$ autour des parités pivots ; banques centrales interviennent pour défendre les monnaies
  \item Chocs pétroliers (1973, 1979): prix du pétrole importé $\uparrow$ $\Rightarrow$ coûts de production $\uparrow$, inflation $\uparrow$, croissance et emploi $\downarrow$
\end{enumerate}

\textbf{Stagflation}: stagnation de la croissance + inflation élevée + chômage persistant.
\begin{enumerate}[\roman*.]
  \item Choc pétrolier $\Rightarrow$ coûts des entreprises $\uparrow$ $\Rightarrow$ prix $\uparrow$ et licenciements
  \item Salaires indexés sur les prix $\Rightarrow$ salaires $\uparrow$ $\Rightarrow$ coûts $\uparrow$ $\Rightarrow$ nouvelle inflation (spirale prix-salaires)
  \item France: inflation $>10\%$ dès 1974, croissance fortement ralentie
\end{enumerate}

\textbf{Alternance relance/rigueur}:
\begin{enumerate}[\roman*.]
  \item 1974 -- Rigueur: encadrement du crédit, taux d'intérêt $\uparrow$, impôts $\uparrow$ $\Rightarrow$ récession et chômage
  \item 1975 -- Relance Chirac: investissements et dépenses publiques $\uparrow$ $\Rightarrow$ reprise temporaire, mais inflation et chômage persistants
  \item 1976--1979 -- Austérité Barre: gel des prix/salaires, prélèvements obligatoires $\uparrow$, crédit limité $\Rightarrow$ inflation encore élevée, investissement faible
  \item 1981 -- Relance Mauroy: SMIC, prestations sociales et emplois publics $\uparrow$ $\Rightarrow$ consommation $\uparrow$ ; croissance en 1982 mais surtout importations et services favorisés
\end{enumerate}

\colbreak

\textbf{Échec des relances keynésiennes}: en situation d'inflation galopante et de dépendance aux importations:
\begin{enumerate}[\roman*.]
  \item Relance budgétaire \\$\Rightarrow$ demande globale $\uparrow$ \\$\Rightarrow$ importations et prix $\uparrow$ \\(déficits jumeaux: public et extérieur)
  \item Inflation française $>$ partenaires \\$\Rightarrow$ compétitivité $\downarrow$ \\$\Rightarrow$ déficit commercial $\uparrow$
  \item Déficit commercial et fuite des capitaux \\$\Rightarrow$ demande de francs $\downarrow$ \\$\Rightarrow$ pression à la dévaluation
  \item Dévaluation \\$\Rightarrow$ importations et pétrole plus chers \\$\Rightarrow$ inflation importée $\uparrow$
\end{enumerate}
Cercle: relance $\to$ inflation $\to$ dévaluation $\to$ inflation.

\textbf{Dévaluation}: baisse officielle de la parité d'une monnaie en changes fixes ; exportations moins chères, mais importations renchéries.

\textbf{Attaques spéculatives}: anticipation d'une dévaluation $\Rightarrow$ vente massive de francs \\$\Rightarrow$ Banque de France vend des devises pour racheter francs \\$\Rightarrow$ réserves de change $\downarrow$ ; si insuffisantes, dévaluation

\textbf{Taux d'intérêt réel}:
\[
\boxed{r_{\text{réel}}\approx r_{\text{nominal}}-\pi}
\]
Inflation élevée $\Rightarrow$ taux réels parfois négatifs malgré des taux nominaux élevés ; politique monétaire peu restrictive.

\textbf{Triangle d'incompatibilité de Mundell-Fleming}: impossible de maintenir simultanément:
\begin{enumerate}[\roman*.]
  \item Fixité des taux de change
  \item Mobilité parfaite des capitaux
  \item Autonomie de la politique monétaire
\end{enumerate}
Avec changes fixes et capitaux libres, un pays ne peut mener durablement une politique monétaire indépendante.

\colbreak
\textbf{France vs Allemagne}:
\begin{enumerate}[\roman*.]
  \item France: inflation élevée, politiques alternées, dévaluations et attaques spéculatives
  \item Allemagne: prioritise stabilité des prix, taux réels positifs, mark fort et crédibilité monétaire
\end{enumerate}

\textbf{Effet d'éviction}: dépenses publiques financées par emprunt $\Rightarrow$ demande de financement $\uparrow$ $\Rightarrow$ taux d'intérêt $\uparrow$ $\Rightarrow$ investissement privé $\downarrow$ 

\subsection{Rigueur (1983--1998)}

\textbf{Tournant de la rigueur (1983)}: abandon de la relance au profit de la désinflation et de la stabilité monétaire pour rester dans le SME (plan Delors).

\textbf{Politique du franc fort / désinflation compétitive}:
\begin{enumerate}[\roman*.]
  \item Rigueur budgétaire: dépenses publiques maîtrisées, impôts $\uparrow$ $\Rightarrow$ déficits et demande limités
  \item Rigueur monétaire: taux d'intérêt élevés pour défendre le franc et lutter contre l'inflation
  \item Désindexation de salaires: rompt spirale prix-salaire
  \item Franc fort: importations moins chères $\Rightarrow$ désinflation importée ; maîtrise des prix/coûts pour retrouver de la compétitivité
  \item Contre-choc pétrolier (1986): pétrole moins cher $\Rightarrow$ désinflation renforcée
\end{enumerate}

\textbf{Crise du SME (1992--1993)}:
\begin{enumerate}[\roman*.]
  \item Réunification allemande (1990) $\Rightarrow$ dépenses publiques et inflation $\uparrow$ $\Rightarrow$ Bundesbank augmente ses taux
  \item Capitaux attirés vers le mark $\Rightarrow$ autres monnaies sous pression ; pays du SME doivent relever leurs taux pour maintenir les parités
  \item Taux élevés $\Rightarrow$ investissement/croissance $\downarrow$ ; attaques spéculatives contre les monnaies européennes
  \item 1993: défense du franc par la BdF ; élargissement des marges du SME de $\pm2,25\%$ à $\pm15\%$
\end{enumerate}
Bilan: changes fixes entre monnaies nationales difficiles à maintenir $\Rightarrow$ motivation pour une monnaie unique.

\textbf{Traité de Maastricht (1992)}: convergence vers l'Union économique et monétaire (UEM).
\begin{enumerate}[\roman*.]
  \item Stabilité des prix, des changes et des taux d'intérêt à long terme
  \item Déficit public de référence $<3\%$ du PIB et dette publique $<60\%$ du PIB
  \item Art. 123: interdiction du financement direct des États par les banques centrales
  \item Coordination et surveillance des politiques économiques nationales
\end{enumerate}

\textbf{Pacte de stabilité et de croissance (PSC, 1997)}:
\begin{enumerate}[\roman*.]
  \item Volet préventif: surveillance des trajectoires budgétaires et des dépenses publiques
  \item Volet correctif: procédure de déficit excessif (PDE), recommandations et sanctions possibles
  \item Révisions: prise en compte du déficit structurel et des circonstances économiques
\end{enumerate}

\textbf{Bilan de la désinflation compétitive}:
\begin{enumerate}[\roman*.]
  \item Succès: inflation de $13,6\%$ (1980) à $0,7\%$ (1998), crédibilité du franc, déficit public ramené à $2,9\%$ du PIB
  \item Limites: chômage élevé, taux d'intérêt contraignants, perte de compétitivité industrielle et inégalités
  \item UEM: taux de change fixes remplacés par monnaie unique ; perte d'autonomie monétaire nationale
\end{enumerate}

\colbreak
\subsection{Relance (1997--2001)}

\textbf{Nouveau policy mix en zone euro}:
\begin{enumerate}[\roman*.]
  \item Politique monétaire: commune, conduite par la BCE ; priorité à la stabilité des prix
  \item Politique budgétaire: nationale, mais encadrée par les critères de Maastricht et le PSC
\end{enumerate}

\textbf{Tournant expansionniste (Jospin)}:
\begin{enumerate}[\roman*.]
  \item Relance budgétaire: dépenses publiques $\uparrow$, impôts des ménages modestes $\downarrow$, politique de l'emploi (35 heures)
  \item Assouplissement monétaire: taux d'intérêt $\downarrow$ $\Rightarrow$ crédits et investissement favorisés
  \item Complémentarité du policy mix: relance budgétaire + monétaire $\Rightarrow$ multiplicateur renforcé, effet d'éviction limité
\end{enumerate}

\subsection{Crises (2007--2026)}

\textbf{Succession de chocs économiques}:
\begin{enumerate}[\roman*.]
  \item 2007--2008: crise des subprimes (bulle immobilière américaine) $\Rightarrow$ faillites bancaires, crise financière mondiale et récession
  \item 2008--2012: crise des dettes souveraines européennes $\Rightarrow$ défiance des créanciers et taux d'emprunt élevés
  \item 2015--2020: risque de déflation $\Rightarrow$ politique monétaire ultra-accommodante
  \item 2020: crise Covid $\Rightarrow$ arrêt partiel de l'activité, aides publiques massives
  \item 2022: guerre en Ukraine et choc énergétique $\Rightarrow$ inflation $\uparrow$, taux directeurs ensuite relevés
\end{enumerate}

\textbf{Crise des dettes souveraines}: déficits et dettes $\uparrow$ après 2008 $\Rightarrow$ défiance des marchés $\Rightarrow$ taux souverains $\uparrow$ $\Rightarrow$ charge d'intérêts $\uparrow$.
\begin{enumerate}[\roman*.]
  \item Ex. Grèce: financement devenu trop coûteux $\Rightarrow$ plans d'aide conditionnés à des mesures de rigueur
  \item FESF puis MES: mécanismes européens de soutien aux États en difficulté pour éviter la contagion
\end{enumerate}

\colbreak
\textbf{Risque de déflation (2015--2020)}:
\begin{enumerate}[\roman*.]
  \item Canal interbancaire bloqué: banques refusent de se prêter malgré des taux bas
  \item Canal du crédit bloqué: banques conservent les liquidités plutôt que de prêter
  \item Faible demande, salaires peu dynamiques, prix de l'énergie bas $\Rightarrow$ inflation sous la cible BCE de $2\%$
\end{enumerate}

\textbf{Politique monétaire non conventionnelle (QE)}: rachats massifs et programmés d'obligations publiques/privées par la BCE sur les marchés.
\begin{enumerate}[\roman*.]
  \item Achats de titres $\uparrow$ $\Rightarrow$ prix des obligations $\uparrow$ \\$\Rightarrow$ rendements/taux longs $\downarrow$
  \item Financement moins coûteux $\Rightarrow$ investissement et croissance soutenus ; dette publique facilitée
  \item APP (2015--2022), PEPP (2020--2022): programmes d'achats d'actifs
  \item Risques: inflation, bulles financières, endettement excessif ; sortie du QE $\Rightarrow$ taux et charge de la dette $\uparrow$
\end{enumerate}

\textbf{Mesures de relance et de soutien (2007--2025)}:
\begin{enumerate}[\roman*.]
  \item Court: sauvetage bancaire (2008), aides Covid, boucliers tarifaires (2022--2023), allègements fiscaux
  \item Moyen/long: Grand Emprunt (2010), France Relance (2021--2022), France 2030 (innovation, numérique, transition écologique)
\end{enumerate}

\textbf{Retour à la rigueur}:
\begin{enumerate}[\roman*.]
  \item Réduction des dépenses: gel des budgets, réformes des retraites (2010, 2023) et de l'assurance-chômage (2019)
  \item Hausse des recettes: impôts sur revenus/capital, réduction des niches fiscales, lutte contre la fraude
  \item Bilan: déficits persistants et dette publique croissante malgré les ajustements
\end{enumerate}

\textbf{Bilan 2026}: France sous procédure de déficit excessif depuis 2024 ; objectif de retour sous les $3\%$ du PIB d'ici 2029. Crédibilité des politiques budgétaires et monétaires essentielle pour conserver la confiance des investisseurs.

\colbreak
\section{Dette publique}

\subsection{Soutenabilité}

\textbf{Solvabilité}: capacité financière d'un État à rembourser sa dette (notion statique). L'État a une durée de vie illimitée, peut lever des impôts et faire rouler sa dette (réemprunter pour rembourser les échéances).

\textbf{Insolvabilité / défaut}: incapacité à honorer la dette $\Rightarrow$ confiance des créanciers $\downarrow$ $\Rightarrow$ taux d'emprunt $\uparrow$ $\Rightarrow$ risque de cercle vicieux.
\begin{enumerate}[\roman*.]
  \item Restructuration: renégocier les échéances, taux ou montants dus
  \item Conséquences: accès aux marchés limité, primes de risque élevées, perte de crédibilité
\end{enumerate}

\textbf{Soutenabilité}: capacité d'un pays à assurer le service de sa dette sans compromettre sa croissance à long terme (notion dynamique). Dépend de la trajectoire dette/PIB et des conditions de financement.

\textbf{Contrainte budgétaire publique}:
\[
\boxed{D_{t+1}=D_t(1+r)-SP_{t+1}}
\]
$D_t$: stock de dette ; $r$: taux d'intérêt nominal apparent (moyen sur la dette) ; $SP$: solde primaire (recettes $-$ dépenses hors intérêts).

\textbf{Ratio dette/PIB}: avec $g$ le taux de croissance nominal du PIB:
\[
\boxed{d_{t+1}=d_t\frac{1+r}{1+g}-sp_{t+1}}
\]
$d=D/PIB$ ; $sp=SP/PIB$.

\textbf{Variation du ratio} (approximation):
\[
\boxed{\Delta d\simeq d_t(r-g)-sp_{t+1}}
\]
\begin{enumerate}[\roman*.]
  \item $r>g$: effet boule de neige ; dette augmente plus vite que le PIB, même si $sp=0$
  \item $g>r$: croissance réduit le poids relatif de la dette ; déficit primaire modéré parfois soutenable
  \item $sp>0$: excédent primaire réduit le ratio ; $sp<0$: déficit primaire l'augmente
\end{enumerate}

\textbf{Solde primaire stabilisant (SPS)}: solde nécessaire pour maintenir le ratio dette/PIB constant ($\Delta d=0$).
\[
\boxed{SPS\simeq d_t(r-g)}
\]
\begin{enumerate}[\roman*.]
  \item $sp>SPS$: endettement/PIB $\downarrow$
  \item $sp=SPS$: ratio stable
  \item $sp<SPS$: endettement/PIB $\uparrow$
\end{enumerate}

\textbf{Limite}: stabiliser une dette ne signifie pas qu'elle soit soutenable ; l'excédent nécessaire peut être économiquement ou politiquement impossible.

\subsection{Stabilisation}

\textbf{Cadre juridique européen}:
\begin{enumerate}[\roman*.]
  \item Maastricht: déficit de référence $3\%$ du PIB, dette $60\%$ du PIB
  \item TSCG (2012): règle d'or limitant le déficit structurel à $0,5\%$ du PIB en principe
  \item France: lois de programmation des finances publiques (LPFP), contrôle par le Haut Conseil des finances publiques (HCFP)
  \item Procédure de déficit excessif (PDE): trajectoire imposée pour réduire déficits et dette, sanctions possibles
\end{enumerate}

\textbf{Situation française}: endettement parmi les plus élevés de la zone euro ; PDE depuis juillet 2024, objectif de déficit inférieur à $3\%$ du PIB en 2029.

\textbf{Dette intergénérationnelle}:
\begin{enumerate}[\roman*.]
  \item Fardeau si dette finance des dépenses courantes non productives ou des déficits structurels
  \item Pas nécessairement un fardeau si elle finance des investissements productifs (éducation, santé, infrastructures, recherche) qui soutiennent la croissance future
\end{enumerate}

\colbreak
\textbf{Risques d'une dette élevée}:
\begin{enumerate}[\roman*.]
  \item Effet boule de neige si $r>g$: poids de la dette $\uparrow$
  \item Debt overhang: dette élevée limite la capacité à intervenir lors d'une nouvelle crise
  \item Perte de confiance: prime de risque et coût d'emprunt $\uparrow$
  \item Perte de souveraineté: aides internationales (ex. FMI) conditionnées à des mesures d'austérité
  \item Consolidation trop brutale: croissance $\downarrow$ $\Rightarrow$ réduction du ratio dette/PIB plus difficile
\end{enumerate}

\textbf{Difficultés d'un excédent primaire}: recettes volatiles, dépenses rigides (santé, retraites), vieillissement, contraintes politiques et électorales.

\subsection{Réduction du déficit}

\textbf{Annuler la dette détenue par l'Eurosystème ?}: la BCE/BdF détient des obligations souveraines achetées sur le marché secondaire (QE). Annulation $\Rightarrow$ baisse comptable du stock de dette.
\begin{enumerate}[\roman*.]
  \item Limite juridique: art. 123 interdit le financement monétaire direct des États
  \item Risque de contagion: autres pays demandent aussi l'annulation $\Rightarrow$ crédibilité de la BCE et de l'euro menacée
  \item Confiance des créanciers $\downarrow$ $\Rightarrow$ taux exigés pour financer l'État $\uparrow$
  \item Dette renouvelée régulièrement: annulation ne résout pas les déficits structurels futurs
\end{enumerate}

\textbf{Consolidation budgétaire}: construire une trajectoire crédible de réduction du déficit sans freiner excessivement croissance et emploi.

\textbf{Étude CAE (octobre 2025)}: scénario central de consolidation sur 6 ans ; effort total estimé à 112 Md€, dont 27 Md€ la première année.
\begin{enumerate}[\roman*.]
  \item Leviers identifiés: économies sur dépenses (108 Md€), recettes supplémentaires (111 Md€), réformes structurelles (45 Md€)
  \item Plusieurs combinaisons possibles ; arbitrages économiques et sociaux
\end{enumerate}

\textbf{Déficit structurel vs conjoncturel}:
\begin{enumerate}[\roman*.]
  \item Conjoncturel: lié aux cycles économiques ; recettes fiscales $\downarrow$ et prestations $\uparrow$ en récession
  \item Structurel: persiste indépendamment de la conjoncture ; déséquilibre durable entre recettes et dépenses
\end{enumerate}
France: environ 85--90\% du déficit primaire attribué au structurel (estimation du cours).

\textbf{Augmenter les recettes}:
\begin{enumerate}[\roman*.]
  \item Impôts sur revenus et patrimoines, réduction des niches fiscales (dépenses fiscales)
  \item Taxe Zucman: proposition d'impôt plancher de $2\%$ sur patrimoines $>100$ M€ ; rendement incertain car réactions des contribuables
\end{enumerate}

\textbf{Effets d'une variation d'impôt}:
\begin{enumerate}[\roman*.]
  \item Effet mécanique: recettes à assiette fiscale constante
  \item Effet comportemental: optimisation, évasion, modification des investissements ou exil fiscal $\Rightarrow$ assiette réduite
\end{enumerate}
Recettes réelles dépendent donc du taux d'imposition et des réactions des contribuables.

\textbf{Réduire les dépenses}:
\begin{enumerate}[\roman*.]
  \item Dépenses de fonctionnement, subventions et aides dont les effets sont faibles
  \item Santé: meilleure coordination des soins, réduction des actes inutiles et des inefficacités
  \item Retraites: vieillissement $\Rightarrow$ dépenses $\uparrow$ ; agir sur âge de départ, emploi des seniors et financement
  \item Préserver les investissements productifs qui soutiennent la croissance potentielle
\end{enumerate}

\textbf{Efficacité des dépenses}: objectif de meilleurs résultats à budget comparable, plutôt qu'une baisse uniforme des dépenses publiques.

\textbf{Fraudes fiscales et sociales}: renforcer les contrôles et le recouvrement ; utilisation de l'IA/data mining pour identifier les irrégularités et améliorer le solde public.

\subsection{Réformes structurelles}

\textbf{Leviers structurels}: améliorer durablement les recettes et réduire la vulnérabilité aux chocs, sans nécessairement produire d'économies budgétaires immédiates.

\textbf{Croissance potentielle}: croissance maximale soutenable avec les capacités de production disponibles, sans tensions inflationnistes.
\begin{enumerate}[\roman*.]
  \item Travail: taux d'emploi $\uparrow$, formation et niveau d'éducation
  \item Capital: investissements, équipements modernes, R\&D
  \item Productivité: progrès technique et meilleure organisation de la production
\end{enumerate}

\textbf{Productivité}:
\begin{enumerate}[\roman*.]
  \item Productivité apparente du travail: valeur ajoutée par heure travaillée
  \item Productivité globale des facteurs (PGF): part de croissance liée à l'efficacité de la combinaison capital/travail (innovation, organisation)
\end{enumerate}
\[
\boxed{P_{\text{travail}}=\frac{VA}{\text{Heures travaillées}}}
\]

\textbf{Gains de productivité}:
\begin{enumerate}[\roman*.]
  \item Coûts de production $\downarrow$ $\Rightarrow$ compétitivité et pouvoir d'achat $\uparrow$
  \item Salaires/production $\uparrow$ $\Rightarrow$ consommation, PIB et recettes fiscales potentiellement $\uparrow$
  \item Destruction de certains emplois à court terme, créations possibles dans d'autres secteurs à long terme (théorie du déversement)
\end{enumerate}

\textbf{Croissance et chômage (loi d'Okun)}: le chômage diminue si la croissance effective dépasse suffisamment la croissance potentielle.
\[
\boxed{g_Y>g_{\text{productivité}}+g_{\text{actifs}}}
\]
Ex.: productivité $+1,3\%$, population active $+1,7\%$ $\Rightarrow$ croissance d'environ $3\%$ pour stabiliser le taux de chômage ; au-delà, chômage peut baisser.

\textbf{Réindustrialisation}: renforcer la production nationale et la souveraineté ; modernisation industrielle, innovation et investissements dans les secteurs stratégiques.

\textbf{Décarbonation}: limiter les émissions et la dépendance aux importations énergétiques ; mécanisme d'ajustement carbone aux frontières (MACF), transition des industries.

\textbf{Bilan}: stabiliser la dette implique maîtrise des dépenses, recettes suffisantes et croissance potentielle renforcée ; nécessité de solidarité, crédibilité et stabilité des politiques économiques.

\end{multicols*}
%%%%%%%%%%%%%%%%%%%%%%%%%%%%%%%%%%%%%%%%%%%%%%%%%%%%%%
%                       End                          %
%%%%%%%%%%%%%%%%%%%%%%%%%%%%%%%%%%%%%%%%%%%%%%%%%%%%%%
\end{document}
